\section{Le robot est-il au bord du cadre$\:$?}
Nous avons vu dans le chapitre précédent que les expressions booléennes {\idt tn}, {\idt tne}, {\idt te}, {\idt tse}, {\idt ts},
 {\idt tsw}, {\idt tw}, {\idt tnw} peuvent être utilisées dans les instructions
 de test ``if'' pour connaître l'orientation du \robot. Une autre  expression booléenne,  
 {\idt tb}, permet de \underline{t}ester si
 le \robot\ est au \underline{b}ord du cadre. Plus précisément {\idt tb} vaut {\idt true} (vrai) 
 si l'exécution supposée de la  commande {\idt av} fait sortir
 le \robot\ de son cadre de déplacement.  L'expression booléenne
 {\idt tb} vaut {\idt false} (faux) si l'exécution supposée 
 de la  commande {\idt av} ne déplace le \robot\ qu'à l'intérieur de son cadre de déplacement.
Par exemple, on peut
 définir une commande {\idt AvSansSortir} qui ordonne au \robot\ d'exécuter un {\idt av} si cela ne le
 fait pas sortir du cadre~:
 
{
\parindent \parindentprog
\medskip

\pagebreak[3] 
\vspace{\neglskip}\mbox{\blc  \blk\idt  \mcl procedure\idt \ AvSansSortir$;$}\par 
\vspace{\neglskip}\mbox{\blc  \blk\idt  \mcl begin\idt }\par 
\pagebreak[3] 
\vspace{\neglskip}\mbox{\blc  \blk\blk\idt  \mcl if\idt \ $($\mcl not\idt \ tb$)$\ \mcl then\idt }\par 
\vspace{\neglskip}\mbox{\blc  \blk\blk\blk\idt  \mcl begin\idt }\par 
\vspace{\neglskip}\mbox{\blc  \blk\blk\blk\blk\idt  av$;$\ }\par 
\vspace{\neglskip}\mbox{\blc  \blk\blk\blk\idt  \mcl end\idt $;$}\par 
\vspace{\neglskip}\mbox{\blc  \blk\idt  \mcl end\idt $;$}\par 

\medskip }

%\pagebreak
\section{Boucle ``while''}

 Une boucle ``while'' permet de répéter l'exécution d'une séquence d'ins\-truc\-tions zéro
ou plusieurs fois en testant à chaque itération s'il faut s'arrêter ou continuer. Le mot anglais
``while'' qui se prononce très approximativement `oua$\cdots$i$\cdots$le' signifie `tant que'. La boucle ``while'' permet de
répéter l'exécution d'une séquence d'instructions tant que la valeur d'une expression booléenne
est {\idt true} (vrai). Par exemple, le programme suivant permet de tracer un trait horizontal au milieu
de l'écran en orientant le \robot\ vers l'est et en répétant des {\idt av} tant que le \robot\ ne sort
pas du cadre~:

 \input TraitHorizontal.tex

Le résultat de l'exécution de ce programme est présenté en réduction dans la figure ci-dessous,
les cadres de déplacement du \robot\ et d'affichage des messages d'erreurs étant en pointillé~:

 \vspace{0.4cm}
 
 \includevspacePostScript[4.4cm]{TraitHorizontal}{0.95}{0.95}

 \vspace{0.4cm}
 
\noindent  Une boucle ``while'' a toujours la forme suivante~:

\index{Boucle ``while''}\index{Expression booléenne}
\begin{center}
\begin{tabular}{@{}|p{12.5cm}|@{}} \hline
{\em Boucle``while''~:}\\ \hline \hline
\begin{minipage}[t]{12.5cm}
{\medskip
\mbox{\blc\ \ \ {\mcl while}\hspace{0.3em}\blanc\hspace{0.3em} {\small\rm Expression booléenne} \hspace{0.3em}\blanc\hspace{0.3em}{\mcl do}}\par 
\mbox{\blc\ \ \ \ \mcl begin}\par 
\mbox{\blc\ \ \ \ \ \small\rm  \ldots\ Instructions répétées (séparées par des `$;$') \ldots}\par 
\mbox{\blc\ \ \ \ {\mcl end}\idt$;$}\par 
\medskip}
\end{minipage}\\ \hline 
\end{tabular}
\end{center}

\noindent La  séquence d'instructions séparées par des `$;$' qui est répétée tant que la valeur de
l'{Expression booléenne} vaut {\idt true} (vrai) s'appelle le {\em corps de
la boucle}\index{Corps de boucle}. Si l'{Expression booléenne}  vaut {\idt false} quand on commence l'exécution de la
boucle ``while'' alors cette exécution se termine immédiatement sans que le corps de la boucle ne soit
exécuté. Si l'{Expression booléenne}  vaut toujours {\idt true} à chaque itération
alors l'exécution de la boucle ``while'' ne se termine jamais$\:$! 
On peut tout de même arrêter l'ordinateur en coupant le courant. Des méthodes plus re\-com\-man\-da\-bles sont expliquées
ci-après. 

\underline{Une erreur grave est de mettre un point-virgule 
`{\idt $;$}' après le {\mcl do}}. Si l'Ex\-pres\-sion booléenne vaut {\idt true}, l'ordinateur boucle éternellement sans 
rien faire. Si l'Expression booléenne vaut {\idt false}, le corps de la boucle est exécuté une fois.

On utilise de préférence les boucles ``for'' (dont l'exécution se termine toujours) en réservant
l'usage des boucles ``while'' au cas où le nombre d'ité\-ra\-tions n'est pas calculable par une expression entière avant de 
commencer à exécuter le corps de la boucle.

\begin{exemple}[Cadre]{\rm Le programme ci-dessous permet de tracer un ca\-dre au\-tour d'un dessin~:

 {\raggedright 

 \vspace{0.4cm}
 
 \includevspacePostScript[4.4cm]{TracerCadre}{0.95}{0.95}
 
 \vspace{0.4cm}
 
\input TracerCadre.tex
 }
\noindent $\Box$ }\end{exemple}

Tous les programmes que nous avions écrits jusqu'alors se terminaient. Ceci veut dire que
 toute exécution du programme prend un temps fini, qui semble parfois long, mais on peut être sûr qu'en attendant suffisamment 
 longtemps, on pourra toujours
 observer un résultat. Si par exemple on a levé le crayon du \robot\ et oublié  de le baisser
 par la suite, le résultat est une feuille blanche mais l'exécution du programme se termine tout de
 même.
 
 Avec la boucle ``while'', on peut écrire des programmes dont l'exécution ne se termine pas.  C'est
 le cas pour le programme suivant~:
 
 \input Boucle1.tex
 
 \noindent Le \robot\ tourne au centre de l'écran, tant que l'ordinateur marche et que l'on n'intervient pas.
 On peut l'arrêter en tapant deux fois au moins sur l'une des touches
 \begin{picture}(12,12)(0, 2) %touche d'echappement
  \put(0,0){\framebox(12,12)[lb]{}}
  \put(1.5,2){\makebox(0,0)[lb]{\tiny esc}}
  \end{picture}\hspace{0.3ex}, 
  \begin{picture}(12,12)(0, 2) %touche §
  \put(0,0){\framebox(12,12)[lb]{}}
  \put(1.5,2){\makebox(0,0)[lb]{\tiny \S}}
  \end{picture}\hspace{0.3ex} ou
  \begin{picture}(12,12)(0, 2) %touche #
  \put(0,0){\framebox(12,12)[lb]{}}
  \put(1.5,2){\makebox(0,0)[lb]{\tiny \#}}
  \end{picture}\hspace{0.3ex}.  L'exécution du programme prend fin car avant 
d'exécuter une commande, le \robot\ regarde si l'utilisateur n'a pas tapé sur l'une de 
ces touches pour lui demander de s'arrêter.
  
  Par contre si l'on exécute le programme \pascal\ suivant~:
  
 \input Boucle2.tex
 
 \noindent l'exécution boucle sans s'arrêter même si l'on tape sur les touches d'arrêt
  \begin{picture}(12,12)(0, 2) %touche d'echappement
  \put(0,0){\framebox(12,12)[lb]{}}
  \put(1.5,2){\makebox(0,0)[lb]{\tiny esc}}
  \end{picture}\hspace{0.3ex}, 
  \begin{picture}(12,12)(0, 2) %touche §
  \put(0,0){\framebox(12,12)[lb]{}}
  \put(1.5,2){\makebox(0,0)[lb]{\tiny \S}}
  \end{picture}\hspace{0.3ex} ou
  \begin{picture}(12,12)(0, 2) %touche #
  \put(0,0){\framebox(12,12)[lb]{}}
  \put(1.5,2){\makebox(0,0)[lb]{\tiny \#}}
  \end{picture}\hspace{0.3ex}
puisque le \robot\ ne reçoit
  aucune commande pendant que l'ordinateur exécute ce programme. 
  La façon d'arrêter l'exécution de ce programme dépend de l'ordinateur utilisé.
  
  Sur un Macintosh, on peut redonner le contrôle au compilateur (si le programme a été
  compilé en mémoire) ou au Finder (si le programme a été
  compilé sur disque) en appuyant sur la touche d'interruption (marquée d'un cercle)
  placée sur le côté du Macintosh~:
 
 \vspace{0.4cm}
 
  \includevspacePostScript[2.9cm]{Interrupt}{0.95}{0.95}

 \vspace{0.4cm}
 
\noindent  Il faut éviter d'appuyer sur la touche de réinitialisation (marquée d'un triangle)
qui revient à éteindre puis à rallumer le Macintosh.
  
 Sur un compatible IBM PC, on peut redonner le contrôle au compilateur (si le programme a été
compilé en mémoire) ou au système MS-DOS (si le programme a été
compilé sur disque) en maintenant la touche contrôle
\begin{picture}(22,12)(0, 2) %touche Ctrl
\put(0,0){\framebox(22,12)[lb]{}}
\put(1.5,2){\makebox(0,0)[lb]{\tiny Ctrl}}
\end{picture}\hspace{0.3ex}
enfoncée et en tapant sur la touche   
\begin{picture}(22,12)(0, 2) %touche Break
\put(0,0){\framebox(22,12)[lb]{}}
\put(1.5,2){\makebox(0,0)[lb]{\tiny Break}}
\end{picture}\hspace{0.3ex}. On peut également, mais c'est déconseillé, réinitialiser l'ordinateur
en maintenant les touches
 \begin{picture}(22,12)(0, 2) %touche Ctrl
\put(0,0){\framebox(22,12)[lb]{}}
\put(1.5,2){\makebox(0,0)[lb]{\tiny Ctrl}}
\end{picture}\hspace{0.3ex}
et
\begin{picture}(14,12)(0, 2) %touche Alt
\put(0,0){\framebox(14,12)[lb]{}}
\put(1.5,2){\makebox(0,0)[lb]{\tiny Alt}}
\end{picture}\hspace{0.3ex}
enfoncées puis en tapant sur la touche
\begin{picture}(14,12)(0, 2) %touche Del
\put(0,0){\framebox(14,12)[lb]{}}
\put(1.5,2){\makebox(0,0)[lb]{\tiny Del}}
\end{picture}\hspace{0.3ex}.

\clearpage
\noindent \rule{137mm}{0.1mm} %---------------------------------------------
\begin{exercices}{Écrire des programmes \pascal\  utilisant des boucles ``while'' pour réaliser les 
dessins suivants~:

\tiret  la frise de la figure \ref{fig:FriTri} ci-dessous$\:$;
% (page \pageref{fig:FriTri})$\:$;

\tiret  la grille des points d'arrêt du \robot\ de la figure \ref{fig:DgQuGr} (page \pageref{fig:DgQuGr}),
sans utiliser bien entendu la commande {\idt dg}$\:$;

\tiret  le quadrillage de la figure \ref{fig:DgQuGr} (page \pageref{fig:DgQuGr})$\:$;

\tiret  le grillage de la figure \ref{fig:DgQuGr} (page \pageref{fig:DgQuGr}) obtenu en superposant deux
qua\-dril\-la\-ges dont un est tourné d'un huitième de tour.

Dans tous les cas le dessin doit être le plus grand possible et s'inscrire à l'intérieur du cadre fixé pour les
déplacements du \robot\ sur l'écran de l'ordinateur.

\noindent $\Box$}\end{exercices}

\begin{exercice}{Définir les commandes {\idt ag},  {\idt ad},  {\idt eh} et  {\idt eb} par des procédures comportant
des boucles ``while'' (on utilisera l'expression booléenne {\idt valBc} qui est vraie si et
seulement si le crayon du \robot\ est baissé).
\noindent $\Box$}\end{exercice}

\begin{figure}[htbp]
 \includevspacePostScript[4.4cm]{FriseDeTriangles}{0.95}{0.95}
 \caption{Frise de triangles (dans un écran réduit)}\label{fig:FriTri}
\end{figure}
\begin{figure}[htbp]
 \includevspacePostScript[14.6cm]{DgQuadGrillage}{0.95}{0.95}
 \caption{Points d'arrêt, quadrillage et grillage sur un écran réduit}\label{fig:DgQuGr}
\end{figure}

%%%%%%%%%%%%%%%%%%%%%%%%%%%%%%%%%%%%%%%%%%%%%%%%%%%%%%%%%
\begin{corrige}{
Nous donnons un programme, les autres se trouvant sur disquette~:
\input FriseDeTriangles.tex
\noindent $\Box$}\end{corrige}
\begin{corrige}{
Nous donnons le corrigé pour la commande {\idt ag}, les autres se trouvant sur disquette~:
\input AgAdEhEb.tex
\noindent $\Box$}\end{corrige}

%%%%%%%%%%%%%%%%%%%%%%%%%%%%%%%%%%%%%%%%%%%%%%%%%%%%%%%%%%%%%%%%%%%%%%%%%%%
\noindent \framebox[\textwidth]{
\begin{minipage}{\myparboxwidth}
\setlength{\parindent}{2em}
\small \sf
\vspace{\abovedisplayskip}
\section*{Addenda}
 
 Une autre forme de répétition non bornée en \pascal\ est la boucle ``repeat'' (qui signifie répéter en anglais).
 Elle a la forme suivante~:
 
 \index{Boucle ``while''}
\begin{center}
\begin{minipage}{10cm}
{\raggedright
\mcl repeat\\
\blc\ \ \small\rm  \ldots\ Instructions répétées (séparées par des `$;$') \ldots\\
\mcl until\hspace{0.3em}\blanc\hspace{0.3em} \small\rm Expression booléenne$;$\\
}\end{minipage}
\end{center}
\vspace{\abovedisplayskip}
\end{minipage}
}
%%%%%%%%%%%%%%%%%%%%%%%%%%%%%%%%%%%%%%%%%%%%%%%%%%%%%%%%%%%%%%%%%%%%%%%%%%%% 

%%%%%%%%%%%%%%%%%%%%%%%%%%%%%%%%%%%%%%%%%%%%%%%%%%%%%%%%%%%%%%%%%%%%%%%%%%%
\noindent \framebox[\textwidth]{
\begin{minipage}{\myparboxwidth}
\setlength{\parindent}{2em}
\small \sf
\vspace{\abovedisplayskip}
\noindent Le corps de la boucle est répété ({\idt repeat}) jusqu'à ({\idt until}) ce que la valeur de 
 l'{\small\rm Expression booléenne} soit égale à {\idt true}. Il est donc toujours exécuté
 au moins une fois.

\indent On peut exprimer les boucles ``while'' à
 l'aide des boucles ``repeat'' et vice-versa. Une boucle ``repeat''~:
 {
\parindent \parindentprog 
\medskip

\pagebreak[3] 
\mbox{\blc  \blk\mcl repeat}\par 
\mbox{\blc  \blk\blk\small\rm Corps de la boucle$;$}\par 
\mbox{\blc  \blk\mcl until {\small\rm Expression booléenne}$;$}\par 

\medskip }

\noindent peut s'écrire avec une boucle ``while'' comme suit~:
{
\parindent \parindentprog
\medskip

\mbox{\blc  \blk\mcl begin}\par 
\mbox{\blc  \blk\blk\small\rm Corps de la boucle$;$}\par 
\mbox{\blc  \blk\blk\mcl while\ not\ \small\rm Expression booléenne\ do}\par 
\mbox{\blc  \blk\blk\blk\mcl begin}\par 
\mbox{\blc  \blk\blk\blk\blk\small\rm Corps de la boucle$;$}\par 
\mbox{\blc  \blk\blk\blk\idt \mcl end$;$}\par 
\mbox{\blc  \blk\mcl end$;$}\par 

\medskip }

\noindent  De la même façon, une boucle ``while''$\:$;
 
{
\parindent \parindentprog
\medskip

\mbox{\blc  \blk\mcl while\ \small\rm Expression booléenne\ \mcl do}\par 
\mbox{\blc  \blk\blk\mcl begin}\par 
\mbox{\blc  \blk\blk\blk\small\rm Corps de la boucle$;$}\par 
\mbox{\blc  \blk\blk\mcl end$;$}\par 

\medskip }

\noindent peut se réécrire avec une boucle ``repeat'' comme suit~:

{
\parindent \parindentprog
\medskip

\mbox{\blc  \blk\mcl if\ \small\rm Expression booléenne\ then}\par 
\mbox{\blc  \blk\blk \mcl begin}\par 
\mbox{\blc  \blk\blk\blk\mcl repeat}\par 
\mbox{\blc  \blk\blk\blk\blk\small\rm Corps de la boucle$;$}\par 
\mbox{\blc  \blk\blk\blk\mcl until\ not\ \small\rm Expression booléenne$;$}\par 
\mbox{\blc  \blk\blk\mcl end$;$}\par 

\medskip }

L'usage de l'une ou l'autre forme d'itération ne dépend que du fait que le corps de la boucle
peut ne jamais être exécuté (boucle ``while'') ou doit toujours l'être au moins une fois
(boucle ``repeat'').

\vspace{\abovedisplayskip}
\end{minipage}
}
%%%%%%%%%%%%%%%%%%%%%%%%%%%%%%%%%%%%%%%%%%%%%%%%%%%%%%%%%%%%%%%%%%%%%%%%%%%% 


